\documentclass[10pt]{report}
%\usepackage{}
\usepackage[T1]{fontenc}
\usepackage[latin1]{inputenc}
%\usepackage[francais]{babel}
\usepackage{a4}
\usepackage{fullpage}
\usepackage{fancyheadings}
\pagestyle{empty}
\usepackage{epsfig}
\setlength{\headrulewidth}{0pt}
\setlength{\footskip}{0.5cm}
\usepackage{anysize}
\marginsize{2.5cm}{2cm}{2cm}{2cm}
\usepackage{url} \urlstyle{tt}


\begin{document}


\thispagestyle{fancy}

\chead{
\ \\[0.5cm]
\makebox[0pt]{
\begin{tabular}{rcr}
\begin{minipage}{25mm}
\centerline{\psfig{file=LORIA.eps,width=25mm}}
\end{minipage}&
\begin{minipage}{10cm}
\begin{center}
{\large\bf Laboratoire Lorrain de Recherche\\ en  Informatique et ses Applications}\\
\noindent UMR N$^\circ$ 7503, CNRS, INRIA\\ UHP-Nancy 1, INPL , Nancy 2
\end{center}
\end{minipage} &
\begin{minipage}{30mm}
\centerline{\psfig{file=INRIA.eps,height=10mm}}
\end{minipage}%\\\hline
\end{tabular}
}}


\cfoot{\small
%hrule
%vspace{2mm}
\makebox[0pt]{
\begin{tabular}{rrrr}
\begin{minipage}{25mm}
\centerline{\psfig{file=CNRS.eps,width=25mm}}
\end{minipage}&
\begin{minipage}{6.6cm}
LORIA, Universit� Henri Poincar�\\
Campus Scientifique - BP 239\\
 54506 Vand{\oe}uvre-l�s-Nancy Cedex, France\end{minipage} &
\begin{minipage}{3.7cm}
Fax :  +33 3 83 59 30 00 \\
T\'el : +33 3 83 27 83 19\\
\url{http://www.loria.fr}
\end{minipage} &
\begin{minipage}{25mm}
\centerline{\psfig{file=logo_nu,width=40mm}}
\end{minipage}
\end{tabular}
}
\ \\[5cm]
}

\newcommand{\edff}[1]{\mbox{\sf EDF$^{(#1)}$}}
\newcommand{\edf}{\mbox{\sf EDF}}

\ \\[2cm]

\hfill
\begin{tabular}{l}
\smallskip\\
St�phane Genaud\\
D�tach� Charg� de Recherches INRIA\\
Tel : +33 (0)389274992\\
{\tt stephane.genaud@loria.fr}\\[5mm]

Nancy,   \today
\end{tabular}
\ \\[1cm]
\noindent{Au Service des Ressources Humaines\\}

Madame, Monsieur,\\

Dans le cadre des offres \emph{Internship} de l'INRIA, nous avons propos�,
Emmanuel Jeannot et moi-m�me, le sujet 
\emph{Design of Latency-aware Communications for Grid Applications (\#320)},
en accord avec le responsable d'�quipe.\\

Nous avons s�lectionn� le candidat suivant:
\begin{center}
Monsieur \textbf{Antonio \textsc{Grassi Chiva}}, domicili\'e \`a:\\
Pablo de Maria 1247, Montevideo, Uruguay\\
email: agrassic@gmail.com\\
Phone: (598 2) 401 59 13\\
Mobile: (598) 099 31 44 31\\
\end{center}
%Facultad de Ingenier�a, Universidad de la Rep�blica

Nous avons convenu avec le candidat que son stage d�buterait le 29 avril 2008 et se terminerait
le 29 ao�t 2008 (dur�e: quatre mois).\\


L'objectif global du stage consiste � mieux comprendre les communications TCP dans un contexte
o� les latences des diff�rentes connexions sont h�t�rog�nes, et de l�, mieux comprendre le
comportement des op�rations de communication collectives de MPI. 
Ce travail s'ins�re dans une des activit�s de l'�quipe qui cherche � am�liorer la perfomance
des communications collectives de MPI dans des environnements de type Grille de calcul.\\

Le travail comportera une partie exp�rimentale importante, effectu�e sur la plateforme Grid5000.
D'autre part, le travail am�nera probablement des propositions algorithmiques pour 
tirer partie de l'h�t�rog�nit� de la plateforme. Dans l'id�al, des contributions significatives
seraient int�gr�es dans une distribution largement distribu�e comme OpenMPI.\\



Je vous prie, Madame, Monsieur, de croire en mes sentiments respectueux.\\







\hfill{St�phane Genaud}\\[2mm]

%\hspace{12cm}\includegraphics[width=4cm]{signature.ps}

\vspace{.2in}

\end{document}



